\documentclass[12pt]{article}
\usepackage{graphicx}

\textwidth 15.2cm
\textheight 22.0cm
\topmargin  -1.5cm
\oddsidemargin 0.5cm
\evensidemargin 0.5cm
%\pagestyle{empty}

%\partopsep -33pt
%\parsep 0pt
%\topskip -33pt
\parskip 0pt
\parindent 0pt
\setlength{\topsep}{0pt}
\setlength{\partopsep}{0pt}
\setlength{\itemsep}{-30pt}
\setlength{\parsep}{-30pt}

\makeatletter
   \def\@listi{\itemsep0pt\parsep0pt}
\makeatother




%\renewcommand{\baselinestretch}{0.9}
\renewcommand{\topfraction}{0.8}
\renewcommand{\textfraction}{0.2}
\renewcommand{\floatpagefraction}{0.8}

%\settowidth{\leftmargin}{CC}
%\addtolength{\leftmargin}{-30pt}
%\setlength{\leftmargin}{5pt}

\begin{document}

\section*{Optimization of functionals with finite width}

The question is what the finite HFB folding width does to the quality
of a functional. Is it beneficial or harmful? To check this, we have
performed a couple of $\chi^2$ fits with systematically varied folding
width. Two series wee done: first, fits of a Skyrme functional to the
data set of SV-min, and second, fits of a Fayans functional to the
data set of Fy($\Delta r$,HFB).  SV-min like fits were done with a
soft pairing cutoff of 30 MeV and with BCS as well as with
HFB. Fy($\Delta r$,HFB) like fits were done all with HFB and two
cutoffs, 15 and 30 MeV. The figure below shoe the resulting quality
measures $\chi^2$ as function of folding width.
%
\bigskip

\centerline{\includegraphics[width=\linewidth]{fit-width.pdf}}

\bigskip
%
The pattern are similar for both functionals. There is a preference
for moderate widths between 0.5 and 0.9 fm. The effect is weak for
SV-min like fits. But Fy($\Delta r$,HFB) like fits gain considerably
from a finite width and they prefer the larger widths which is rather
welcome. The good news is that width of about 1 fm are beneficial for
Fayans functionals, already from quality alone. Moreover, we gain the
intended better convergence with these widths.

\end{document}
